\documentclass[11pt]{article}

\usepackage[margin=1in]{geometry}
\usepackage[T1]{fontenc}
\usepackage[utf8]{inputenc}
\usepackage{amsmath,amssymb,amsthm,mathtools}
\usepackage{bm}
\usepackage{booktabs}
\usepackage{array}
\usepackage{longtable}
\usepackage{tabularx}
\usepackage{multirow}
\usepackage{xcolor}
\usepackage{enumitem}
\usepackage{hyperref}
\usepackage{float}
\usepackage{caption}
\usepackage{titlesec}
\usepackage{listings}

\hypersetup{
  colorlinks=true,
  linkcolor=blue,
  urlcolor=blue,
  citecolor=blue
}

\titleformat{\section}{\large\bfseries}{\thesection}{0.6em}{}
\titleformat{\subsection}{\normalsize\bfseries}{\thesubsection}{0.6em}{}
\titleformat{\subsubsection}{\normalsize\bfseries}{\thesubsubsection}{0.6em}{}

\lstset{
  basicstyle=\ttfamily\small,
  breaklines=true,
  frame=single,
  backgroundcolor=\color{gray!8},
  commentstyle=\color{gray},
  keywordstyle=\color{blue},
}

\definecolor{noteblue}{RGB}{220,235,252}
\newcommand{\note}[1]{\colorbox{noteblue}{\parbox{\dimexpr\linewidth-2\fboxsep}{\small #1}}}

\title{\textbf{End-to-End Music Recommendation System\\Architecture Specification}}
\author{SpotyBoys MLOps}
\date{2026}

\begin{document}
\maketitle
\tableofcontents
\newpage

% =============================================================================
\section{System Overview}
% =============================================================================

This document specifies the production architecture for a real-time personalised music
recommendation system built on top of Navidrome, an open-source music streaming server.
The system combines three ML components --- Item2Vec (C1), Multi-Recall Retriever (C2), and
GRU Ranker (C3) --- with a curated music library to deliver next-track recommendations
without requiring users to provide any explicit input beyond normal playback interactions.

\subsection{Design Principles}

\begin{itemize}[leftmargin=1.5em]
  \item \textbf{Closed catalog}: the music library is curator-controlled.  Every track served
        to users has a pre-computed embedding in \texttt{item2vec\_128d.npy}, guaranteeing that
        C2 and C3 can score any candidate.
  \item \textbf{Zero user friction}: labels are derived automatically from playback behaviour
        (play ratio), not from explicit ratings.
  \item \textbf{Latency budget}: the full C2+C3 pipeline must return recommendations in
        $\leq 200\,\text{ms}$ at serving time.  Navidrome audio delivery is decoupled from
        recommendation latency.
  \item \textbf{Separation of concerns}: Navidrome handles auth, streaming, and metadata
        display; the Recommendation API handles all ML inference.
\end{itemize}

\subsection{High-Level Component Diagram}

\begin{center}
\begin{tabular}{|c|}
\hline
\\
\textbf{User Browser / Mobile App} \\
\quad(custom lightweight web player) \\
\\
\hline
\end{tabular}
\\[4pt]
$\downarrow$ \textit{audio stream (Subsonic API)} \quad\quad
$\downarrow$ \textit{play events + recommendation requests (REST)}
\\[4pt]
\begin{tabular}{|c|c|}
\hline
\quad\textbf{Navidrome}\quad & \quad\textbf{Recommendation API (FastAPI)}\quad \\
port 4533 & port 8080 \\
\hline
\end{tabular}
\\[4pt]
$\downarrow$ \textit{music files} \qquad\qquad
$\downarrow$ \textit{session / u\_long cache}
\quad $\downarrow$ \textit{interaction logs}
\quad $\downarrow$ \textit{model artifacts}
\\[4pt]
\begin{tabular}{|c|c|c|}
\hline
\quad\textbf{S3}\quad & \quad\textbf{Redis}\quad & \quad\textbf{PostgreSQL}\quad \\
artifacts + audio & in-request cache & persistent store \\
\hline
\end{tabular}
\end{center}

% =============================================================================
\section{Deployment Topology}
% =============================================================================

\subsection{Virtual Machines}

The system uses two VMs:

\begin{longtable}{@{}llp{7cm}@{}}
\toprule
VM & Role & Services \\
\midrule
\textbf{proj23-training-part} (persistent) &
  Serving + storage &
  Navidrome, Recommendation API, Redis, PostgreSQL, MLflow \\[4pt]
\textbf{proj23-training} (GPU, ephemeral) &
  Model training &
  PyTorch training jobs; \texttt{MLFLOW\_TRACKING\_URI} points to persistent VM \\
\bottomrule
\end{longtable}

\subsection{Docker Services on Persistent VM}

All services run as Docker containers managed by \texttt{docker-compose}.

\begin{longtable}{@{}lllp{5.5cm}@{}}
\toprule
Container & Image & Port & Purpose \\
\midrule
\texttt{navidrome} & \texttt{deluan/navidrome:latest} & 4533 &
  Music streaming; Subsonic API for audio delivery and user auth \\[4pt]
\texttt{rec-api} & custom FastAPI image & 8080 &
  ML inference: C2 retrieval + C3 ranking; play event ingestion \\[4pt]
\texttt{redis} & \texttt{redis:7-alpine} & 6379 &
  In-memory session cache and u\_long cache \\[4pt]
\texttt{postgres} & \texttt{postgres:16} & 5432 &
  Interaction logs, user/track mapping tables \\[4pt]
\texttt{mlflow} & custom MLflow image & 8000 &
  Experiment tracking; backend = PostgreSQL; artifacts = S3 \\
\bottomrule
\end{longtable}

\subsection{Persistent Storage}

\begin{itemize}[leftmargin=1.5em]
  \item PostgreSQL data directory: \texttt{/mnt/persist/postgres\_data} (mounted volume)
  \item Navidrome music library: \texttt{/mnt/persist/music} (mounted volume)
  \item Navidrome database: \texttt{/mnt/persist/navidrome\_data} (mounted volume)
\end{itemize}

% =============================================================================
\section{Online Serving Data Flow}
% =============================================================================

This section describes the complete lifecycle of a single user interaction, from a track
finishing playback to the next recommendation being delivered.

\subsection{Label Derivation from Playback Behaviour}

Labels are derived \emph{client-side} using the browser's HTML5 Audio API.
The web player attaches listeners to the \texttt{HTMLAudioElement}:

\begin{itemize}[leftmargin=1.5em]
  \item On \texttt{play} event: record \texttt{t\_start = audio.currentTime} and wall-clock
        timestamp.
  \item On \texttt{ended} event (natural completion): \texttt{play\_duration\_ms $\approx$
        track\_duration\_ms} $\Rightarrow$ always \texttt{positive}.
  \item On user skip / next-track / app close: read \texttt{audio.currentTime} at that moment
        to obtain the actual seconds listened.
\end{itemize}

The computed \texttt{play\_duration\_ms} (how long the user actually listened) and
\texttt{track\_duration\_ms} (total length of the track, from Navidrome metadata) are
sent together in the \texttt{POST /event} body.  The server derives the label from the ratio
$r = \text{play\_duration\_ms} / \text{track\_duration\_ms}$:
\[
  \text{label} =
  \begin{cases}
    \texttt{positive} & r \geq 0.80 \\
    \texttt{skip}     & r \leq 0.20 \\
    \texttt{neutral}  & \text{otherwise}
  \end{cases}
\]

\note{\textbf{Note}: the server does \emph{not} store separate play-start / play-end
timestamps.  The client is responsible for measuring elapsed playback time via
\texttt{audio.currentTime}; the server only stores the pre-computed duration pair
and the derived label.}

\subsection{Step-by-Step Request Lifecycle}

\begin{enumerate}[leftmargin=1.5em]

  \item \textbf{Track finishes (or user skips).}
        The web player reads \texttt{audio.currentTime} at the moment of skip or
        \texttt{ended} event, computes \texttt{play\_duration\_ms}, and POSTs to the
        Recommendation API:
\begin{lstlisting}[language=json]
POST /event
{
  "user_id": 1042,
  "track_id": 58291,
  "session_uuid": "a3f9...",
  "play_duration_ms": 198000,
  "track_duration_ms": 214000
}
\end{lstlisting}

  \item \textbf{Recommendation API derives label} using the ratio formula defined in
        Section~3.1.

  \item \textbf{State updates (async, non-blocking).}
        \begin{itemize}
          \item \textbf{Redis}: append \texttt{(track\_id, label)} to
                \texttt{session:\{user\_id\}}; cap list at 20 entries (matches $L=20$).
          \item \textbf{PostgreSQL}: insert one row into \texttt{interactions}
                (fire-and-forget; does not block the response).
        \end{itemize}

  \item \textbf{Client requests next $N$ recommendations.}
\begin{lstlisting}
GET /recommend?user_id=1042&n=10
\end{lstlisting}

  \item \textbf{Recommendation API assembles context.}
        \begin{itemize}
          \item Read \texttt{session:\{user\_id\}} from Redis
                $\Rightarrow$ \texttt{session\_track\_ids}, \texttt{session\_labels}.
          \item Read \texttt{ulong:\{user\_id\}} from Redis (128D vector, cached).
                If cache miss: compute via nearest centroid from
                \texttt{user\_centroids.pkl}; write back with TTL\,=\,24\,h.
          \item Read \texttt{seen:\{user\_id\}} from Redis (recently served track set).
        \end{itemize}

  \item \textbf{C2 --- Multi-Recall Retrieval.}
        \texttt{MultiRecallRetriever.retrieve(user\_id, session\_track\_ids, session\_labels)}
        returns up to 200 \texttt{(track\_id, score)} pairs via three branches:
        \begin{itemize}
          \item Co-occurrence branch (up to 100): recency-weighted $C_\text{sess}$ + $C_\text{pl}$
          \item Preference-NN branch (up to 80): dot-product against user centroids
          \item Popularity fallback (up to 20): pre-sorted \texttt{pop\_scores.csv}
        \end{itemize}
        Candidates present in \texttt{seen:\{user\_id\}} are filtered out.

  \item \textbf{C3 --- GRU Ranker scoring.}
        \texttt{GRURankerInference.score(user\_id, session\_track\_ids, session\_labels, candidates)}
        scores all candidates in a single batched forward pass and returns them sorted by
        logit descending.

  \item \textbf{Response.}
        Top $N$ \texttt{track\_id}s are returned; the web player adds them to the playback queue
        and fetches their audio streams from Navidrome using the Subsonic
        \texttt{stream?id=\{navidrome\_track\_id\}} endpoint.
        The returned track IDs are appended to \texttt{seen:\{user\_id\}} (TTL\,=\,1\,h).

\end{enumerate}

\subsection{Latency Budget}

\begin{longtable}{@{}lrl@{}}
\toprule
Step & Target & Notes \\
\midrule
Redis session read & $<1\,\text{ms}$ & in-memory \\
C2 retrieval & $<30\,\text{ms}$ & sparse matrix ops; embeddings in RAM \\
C3 ranking (200 candidates) & $<50\,\text{ms}$ & CPU forward pass; batch size 200 \\
PostgreSQL write (async) & --- & does not block response \\
\midrule
\textbf{Total (p99)} & $\mathbf{<150\,\text{ms}}$ & \\
\bottomrule
\end{longtable}

% =============================================================================
\section{Storage Architecture}
% =============================================================================

\subsection{Redis --- In-Request Cache}

Redis holds all state needed to serve a recommendation request with sub-millisecond reads.
No ML artifacts are stored in Redis; only small derived vectors and ID lists.

\begin{longtable}{@{}lllp{4.5cm}@{}}
\toprule
Key pattern & Redis type & Value & TTL \\
\midrule
\texttt{session:\{user\_id\}} & List &
  JSON-encoded \texttt{[(track\_id, label), \ldots]} capped at 20 &
  8 hours (reset on new session) \\[4pt]
\texttt{ulong:\{user\_id\}} & String &
  Base64-encoded float32 array, shape (128,) &
  24 hours \\[4pt]
\texttt{seen:\{user\_id\}} & Set &
  track\_ids served in the last hour &
  1 hour \\
\bottomrule
\end{longtable}

\paragraph{Session boundary.}
A new session UUID is generated by the web player when the app starts or after 30 minutes
of inactivity.  On session start, the existing \texttt{session:\{user\_id\}} key is deleted
and a fresh list is created.

\subsection{PostgreSQL --- Persistent Relational Store}

PostgreSQL stores data that must survive VM restarts and is used for offline retraining.
It is \emph{never} on the critical path of a live recommendation request.

\subsubsection{Schema}

\begin{lstlisting}[language=SQL]
-- User mapping: Navidrome user_id <-> internal integer user_id
CREATE TABLE users (
    user_id          BIGSERIAL PRIMARY KEY,
    navidrome_user_id VARCHAR(64) UNIQUE NOT NULL,
    created_at       TIMESTAMP DEFAULT NOW()
);

-- Track mapping: Navidrome track_id <-> internal integer track_id
-- track_id matches the row index in item2vec_track_to_row.json
CREATE TABLE tracks (
    track_id           INTEGER PRIMARY KEY,
    navidrome_track_id VARCHAR(64) UNIQUE NOT NULL,
    title              VARCHAR(512),
    artist             VARCHAR(256),
    duration_ms        INTEGER
);

-- Raw interaction log (source of truth for retraining)
CREATE TABLE interactions (
    id                BIGSERIAL PRIMARY KEY,
    user_id           BIGINT    NOT NULL REFERENCES users(user_id),
    track_id          INTEGER   NOT NULL REFERENCES tracks(track_id),
    session_uuid      VARCHAR(36) NOT NULL,
    play_duration_ms  INTEGER   NOT NULL,
    track_duration_ms INTEGER   NOT NULL,
    label             VARCHAR(8) NOT NULL,   -- positive / neutral / skip
    created_at        TIMESTAMP DEFAULT NOW()
);

CREATE INDEX idx_interactions_user_time
    ON interactions (user_id, created_at DESC);
CREATE INDEX idx_interactions_created
    ON interactions (created_at DESC);
\end{lstlisting}

\subsubsection{Usage Pattern}

\begin{itemize}[leftmargin=1.5em]
  \item \textbf{Writes}: one row per play event, asynchronously via the Recommendation API.
  \item \textbf{Reads}: batch export only (nightly cron → S3); never read at serving time.
\end{itemize}

\subsection{S3 --- Object Storage}

All large binary artifacts and raw data exports live in S3.  The bucket
\texttt{proj23-mlflow-artifacts} is organised as follows:

\begin{lstlisting}
proj23-mlflow-artifacts/
|
|-- artifacts/
|   |-- item2vec/
|   |   |-- item2vec_128d.npy           # (746609, 128) float32 embedding matrix (365 MB)
|   |   |-- item2vec_track_to_row.json  # track_id -> row index (13 MB)
|   |   `-- item2vec_catalog.csv        # canonical track list
|   |
|   |-- retriever/
|   |   |-- cooc/
|   |   |   |-- cooc_session.npz        # C_sess sparse CSR (8.2M nnz, 26 MB)
|   |   |   `-- cooc_playlist.npz       # C_pl sparse CSR (9.1M nnz)
|   |   |-- pref_nn/
|   |   |   `-- user_centroids.pkl      # dict[user_id -> [(centroid, size)]] (110 MB)
|   |   `-- popularity/
|   |       `-- pop_scores.csv          # 746K tracks sorted by log(1+count)
|   |
|   `-- ranker/
|       |-- gru_ranker.pt               # best model checkpoint (NDCG@5)
|       `-- gru_ranker_config.json      # d_emb, n_layers, dropout, L
|
|-- interactions/
|   |-- year=2026/month=04/day=14/
|   |   `-- interactions_20260414.parquet
|   `-- ...                             # one parquet file per day
|
`-- mlflow/                             # MLflow artifact store (auto-managed)
    `-- <experiment_id>/<run_uuid>/artifacts/
\end{lstlisting}

\paragraph{Artifact loading at startup.}
When the Recommendation API container starts, it downloads the four model artifacts
(\texttt{item2vec\_128d.npy}, \texttt{item2vec\_track\_to\_row.json},
\texttt{user\_centroids.pkl}, \texttt{gru\_ranker.pt}) from S3 into memory.
Total startup RAM cost: $\approx 490\,\text{MB}$.

% =============================================================================
\section{Recommendation API Design}
% =============================================================================

The Recommendation API is a FastAPI service that exposes two endpoints.  It loads
\texttt{MultiRecallRetriever} and \texttt{GRURankerInference} once at startup and
keeps them in process memory.

\subsection{Endpoints}

\subsubsection{\texttt{POST /event} --- Ingest Play Event}

\begin{lstlisting}[language=json]
// Request body
{
  "user_id":          1042,
  "track_id":         58291,
  "session_uuid":     "a3f9c2d1-...",
  "play_duration_ms": 198000,
  "track_duration_ms": 214000
}

// Response 200
{
  "label": "positive"
}
\end{lstlisting}

\textbf{Processing} (all synchronous except PostgreSQL write):
\begin{enumerate}[leftmargin=1.5em, topsep=2pt]
  \item Derive \texttt{label} from play ratio.
  \item \texttt{RPUSH session:\{user\_id\} (track\_id, label)}; trim to last 20 entries.
  \item Invalidate \texttt{ulong:\{user\_id\}} if new label is \texttt{positive}
        (session mean may have shifted meaningfully).
  \item Fire-and-forget INSERT into \texttt{interactions}.
\end{enumerate}

\subsubsection{\texttt{GET /recommend} --- Get Next Tracks}

\begin{lstlisting}
GET /recommend?user_id=1042&n=10
\end{lstlisting}

\begin{lstlisting}[language=json]
// Response 200
{
  "user_id": 1042,
  "recommendations": [
    {"rank": 1, "track_id": 73812, "score": 2.341},
    {"rank": 2, "track_id": 41205, "score": 1.987},
    ...
  ]
}
\end{lstlisting}

\textbf{Processing}:
\begin{enumerate}[leftmargin=1.5em, topsep=2pt]
  \item Read session from Redis.
  \item Resolve \texttt{u\_long} (Redis cache or compute from centroids).
  \item C2: \texttt{MultiRecallRetriever.retrieve(\ldots)} $\to$ up to 200 candidates.
  \item Filter out \texttt{seen:\{user\_id\}}.
  \item C3: \texttt{GRURankerInference.score(\ldots)} $\to$ sorted list.
  \item Return top $n$; add returned IDs to \texttt{SADD seen:\{user\_id\}}.
\end{enumerate}

\subsection{Navidrome $\leftrightarrow$ Track ID Mapping}

Navidrome assigns its own opaque string IDs to tracks.  At library ingestion time, an
admin script populates the \texttt{tracks} table by calling the Navidrome Subsonic API
(\texttt{getMusicDirectory}) and matching each track's file path or MusicBrainz ID to a
row in \texttt{item2vec\_track\_to\_row.json}.

The web player always communicates with the Recommendation API using \emph{internal}
\texttt{track\_id} (integer), and fetches audio streams from Navidrome using
\emph{navidrome\_track\_id} (string), resolved via a local lookup table loaded at startup.

% =============================================================================
\section{Retraining Loop}
% =============================================================================

The retraining loop ensures the model stays aligned with evolving user preferences.
It is fully automated via cron jobs on the persistent VM.

\subsection{Overview}

\begin{center}
\begin{tabular}{ccccc}
PostgreSQL & $\xrightarrow{\text{daily export}}$ & S3 parquet &
$\xrightarrow{\text{weekly build}}$ & new training data \\
& & & $\downarrow$ & \\
& & updated artifacts & $\xleftarrow{\text{retrain}}$ & GPU VM \\
& & $\downarrow$ & & \\
& & Rec API hot-reload & & \\
\end{tabular}
\end{center}

\subsection{Step 1: Daily Interaction Export (cron: 02:00 UTC)}

A cron job on the persistent VM runs:

\begin{lstlisting}[language=bash]
# Export yesterday's interactions from PostgreSQL to S3
DATE=$(date -d "yesterday" +%Y-%m-%d)
docker exec postgres psql -U user -d appdb -c \
  "COPY (SELECT * FROM interactions
         WHERE created_at::date = '${DATE}')
   TO STDOUT WITH CSV HEADER" \
  | aws --endpoint-url ${S3_ENDPOINT} \
        s3 cp - s3://proj23-mlflow-artifacts/interactions/year=.../
\end{lstlisting}

\subsection{Step 2: Weekly Feature Rebuild (cron: Sunday 03:00 UTC)}

On the GPU VM (triggered via SSH or remote cron), the following pipeline runs:

\begin{enumerate}[leftmargin=1.5em]
  \item \textbf{Download new interactions} from S3 and merge with historical data.
  \item \textbf{Rebuild user centroids}:
\begin{lstlisting}[language=bash]
uv run python -m src.retriever.pref_nn.build
\end{lstlisting}
  \item \textbf{Rebuild ranker training data}:
\begin{lstlisting}[language=bash]
uv run python -m src.ranker.data.build
\end{lstlisting}
  \item \textbf{Retrain GRU Ranker}:
\begin{lstlisting}[language=bash]
uv run python -m src.ranker.train \
    --mlflow-experiment gru-ranker-weekly \
    --run-name weekly-$(date +%Y%m%d)
\end{lstlisting}
  \item \textbf{Upload new artifacts} to S3:
\begin{lstlisting}[language=bash]
bash upload_results.sh
\end{lstlisting}
\end{enumerate}

\subsection{Step 3: Model Hot-Reload on Serving VM}

After new artifacts land in S3, the Recommendation API reloads the model without downtime:

\begin{lstlisting}[language=bash]
# On persistent VM
docker exec rec-api python -c "
import requests
requests.post('http://localhost:8080/admin/reload')
"
\end{lstlisting}

The \texttt{/admin/reload} endpoint re-instantiates \texttt{GRURankerInference} and
\texttt{MultiRecallRetriever} from the latest S3 artifacts, then atomically swaps the
references used by the request handlers.

\subsection{Retraining Trigger Thresholds}

\begin{longtable}{@{}llp{6cm}@{}}
\toprule
Trigger & Condition & Action \\
\midrule
Time-based & Every Sunday 03:00 UTC & Full retrain regardless of data volume \\
Data-based & $>50{,}000$ new interactions & Force early retrain \\
Metric-based & NDCG@5 drops $>5\%$ (monitored via MLflow) & Alert + manual retrain \\
\bottomrule
\end{longtable}

% =============================================================================
\section{Cold Start Handling}
% =============================================================================

\subsection{New User (no listening history)}

\begin{itemize}[leftmargin=1.5em]
  \item \texttt{user\_centroids} has no entry for this user.
  \item \texttt{u\_long} is set to $\mathbf{0}_{128}$ (zero vector).
  \item C2 falls back entirely to the popularity branch (20 tracks) and co-occurrence
        seeded from the first played track's co-occurrence row.
  \item After the first 3 positive interactions in a session, the session mean embedding
        provides a meaningful \texttt{u\_long} approximation.
\end{itemize}

\subsection{New Track (not in item2vec catalog)}

Since the music library is curator-controlled and all tracks are ingested before
serving begins, this case should not occur in normal operation.  The admin
ingestion script must verify that every uploaded track has a valid row in
\texttt{item2vec\_track\_to\_row.json} before making it playable.  Tracks that fail
this check are flagged as \texttt{not\_in\_catalog} and excluded from recommendations.

% =============================================================================
\section{Summary of Data Ownership}
% =============================================================================

\begin{longtable}{@{}lll@{}}
\toprule
Data & Store & Lifetime \\
\midrule
Current session context & Redis & 8 hours \\
u\_long vector cache & Redis & 24 hours \\
Recently served tracks & Redis & 1 hour \\
\midrule
User--Navidrome ID mapping & PostgreSQL & Permanent \\
Track--Navidrome ID mapping & PostgreSQL & Permanent \\
Raw interaction logs & PostgreSQL + S3 & Permanent \\
\midrule
Item2Vec embeddings & S3 (+ RAM at runtime) & Until next C1 retrain \\
Co-occurrence matrices & S3 (+ RAM at runtime) & Until next C2 rebuild \\
User centroids & S3 (+ RAM at runtime) & Until next weekly rebuild \\
GRU Ranker weights & S3 (+ RAM at runtime) & Until next weekly retrain \\
MLflow experiment records & PostgreSQL + S3 & Permanent \\
\bottomrule
\end{longtable}

\end{document}
